\section{Defining the Task}
\label{sec:defining-task}

Admmit, a Norwegian software firm building products for the transport sector, offered a bachelor task through NTNU under project code OPG29. The task description asked for a working prototype of a planning platform that could help Norwegian traffic coordinators assign drivers and vehicles to the day's transport assignments. The two-student team accepted the task and committed to delivering both the artefact, named RP, and the thesis that surrounds it. From that point onwards, two workstreams ran in parallel: building the platform and conducting the research that would inform its design. Both students contributed to both workstreams over the course of the project.

The task description fixed the topic but not the empirical foundation, and the team needed answers before any planning interface could be designed responsibly. Two questions in particular: how do Norwegian traffic coordinators actually plan a day's assignments, and what software, if any, do they use today? Neither question could be answered from the desk. The team turned to Admmit's customer roster, called seven traffic coordinators and transport managers directly, and conducted brief semi-structured calls of roughly fifteen minutes each. The calls were spread across the first half of March 2026 rather than compressed into one day; each call was conducted on its own, and the team adjusted prompts and reflected between calls so that gaps in earlier conversations could be probed in later ones. Vendor names surfaced gradually through "what tools do you use" questions: several companies described mixes rather than single-vendor adoption, with Timpex, Opter, and internal tools appearing in different combinations across the seven companies interviewed.

Admmit had already specified two requirements before the first call. The first was Human-in-the-Loop (HITL) operation: the algorithm proposes a plan, and the coordinator can inspect, modify, accept, or reject every assignment in it. The second was a multi-tenant architecture, meaning one deployed system serving several customer companies, with each company's data and configuration isolated from the others. After the calls, the team agreed with what Admmit had set. The seven coordinators voiced a clear preference for a system that suggests a plan they can correct rather than one that decides for them, which matched the HITL requirement. Their accounts of differing operational rules across companies matched the case for the multi-tenant architecture. Admmit had the requirements; the seven companies interviewed validated them.


